% ============================================================================
% Fichier  : partie3/P03_C08_clauses_4_6.tex
% Titre    : Contexte, Leadership, Planification
% Auteur   : MINKA MI NGUIDJOÏ Thierry Emmanuel
% ============================================================================
% Sources : ISO/IEC 27001:2022, clauses 4 à 6. Mairie d'Eséka : cas fictif.
% ============================================================================

\chapter{Contexte, Leadership, Planification}
\label{chap:clauses46}
\chaptermark{Clauses 4 à 6}

\epigraph{%
    « Avant de planifier la sécurité, il faut savoir pour qui,\\
    contre quoi et dans quelles limites. »%
}{Principe du cours}

\niveauChapitre{Intermédiaire}{Le Chapitre~\ref{chap:iso27000}.}

\begin{boxobjectifs}
\begin{itemize}
    \item Déterminer le contexte et les parties intéressées (clause~4).
    \item Préciser le domaine d'application du \sigle{SMSI}.
    \item Caractériser le leadership exigé de la direction (clause~5).
    \item Planifier les actions face aux risques et les objectifs (clause~6).
    \item Comprendre le lien entre objectifs de sécurité et appréciation des
          risques.
\end{itemize}
\end{boxobjectifs}

\bigskip

% ============================================================================
\section*{Question directrice : comment poser les fondations d'un SMSI conforme ?}
% ============================================================================

Les trois premières clauses d'exigences fixent le cadre~: qui, quoi, pourquoi,
et avec quels objectifs. Ce chapitre les traduit pour la mairie d'Eséka.

\medskip
\textbf{Plan du chapitre.}
\begin{enumerate}
    \item Le contexte (clause~4) (section~\ref{sec:c08-contexte}).
    \item Le domaine d'application (section~\ref{sec:c08-domaine}).
    \item Le leadership (clause~5) (section~\ref{sec:c08-leadership}).
    \item La planification (clause~6) (section~\ref{sec:c08-planif}).
    \item Les objectifs de sécurité (section~\ref{sec:c08-objectifs}).
\end{enumerate}

\separateur

% ============================================================================
\section{Le Contexte (Clause 4)}
\label{sec:c08-contexte}
% ============================================================================

\begin{boxdef}
La clause~4 exige de déterminer les \termecle{enjeux internes et externes}
pertinents, les \termecle{parties intéressées} et leurs attentes, puis de
fixer le \termecle{domaine d'application} du \sigle{SMSI}.
\end{boxdef}

\begin{tableaucours}{L{3.2cm}L{9.0cm}}{%
    \tableauHeader{Élément} &
    \tableauHeader{Pour la mairie d'Eséka}
}
Enjeux externes & Loi 2024/017, attentes des usagers, menace cyber \\
Enjeux internes & Budget contraint, compétences limitées, matériel vieillissant \\
Parties intéressées & Usagers, tutelle, ANTIC, prestataires, agents \\
\end{tableaucours}
\vspace{-0.8em}
\begin{center}{\small\itshape Tableau~8.1. Le contexte de la mairie (clause 4)}\end{center}

% ============================================================================
\section{Le Domaine d'Application}
\label{sec:c08-domaine}
% ============================================================================

Le \termecle{domaine d'application} (\emph{scope}) dit ce que couvre le
\sigle{SMSI}~: quels services, sites, actifs et exclusions. Il doit être écrit
et justifié~: un périmètre flou rend l'audit impossible.

\begin{boxprocedure}
\textbf{Constatation.} La mairie fixe le domaine~: « le système d'information
des services aux usagers de la mairie d'Eséka (état civil et taxes), à
l'exclusion du réseau technique de la voirie ».

\medskip
\textbf{Analyse.} L'exclusion est explicite et justifiée (la voirie ne traite
pas de données sensibles). Un domaine clair protège contre l'extension
incontrôlée et cadre l'audit.
\end{boxprocedure}

% ============================================================================
\section{Le Leadership (Clause 5)}
\label{sec:c08-leadership}
% ============================================================================

La clause~5 exige que la direction \termecle{démontre} son leadership~: établir
la politique, attribuer rôles et responsabilités, fournir les ressources,
promouvoir l'amélioration continue. Ce n'est pas une déclaration~: ce sont des
actes vérifiables (Chapitre~\ref{chap:gouvernance}).

\begin{boxmetier}{Direction générale}
\textbf{Des preuves de leadership.} Le maire prouve son engagement par une
politique signée, un budget voté, une RSSI nommée et sa présence au comité.
L'auditeur demandera ces preuves~; les bonnes intentions ne comptent pas.
\end{boxmetier}

% ============================================================================
\section{La Planification (Clause 6)}
\label{sec:c08-planif}
% ============================================================================

\begin{boxdef}
La clause~6 exige de planifier les \termecle{actions face aux risques et
opportunités}, de définir le \termecle{processus d'appréciation et de
traitement des risques}, et de fixer des \termecle{objectifs de sécurité}
mesurables.
\end{boxdef}

La planification est le cœur « Plan » du \sigle{PDCA}. C'est ici que la norme
renvoie à l'appréciation des risques (partie IV)~: on ne planifie pas au
hasard, mais au vu des risques appréciés.

% ============================================================================
\section{Les Objectifs de Sécurité}
\label{sec:c08-objectifs}
% ============================================================================

\begin{boxdef}
Un \termecle{objectif de sécurité} est un résultat visé, cohérent avec la
politique, \emph{mesurable}, suivi et mis à jour. « Améliorer la sécurité »
n'est pas un objectif~; « sauvegarder quotidiennement l'état civil et le
vérifier chaque mois » en est un.
\end{boxdef}

\begin{tableaucours}{L{5.4cm}L{3.0cm}L{3.4cm}}{%
    \tableauHeader{Objectif} &
    \tableauHeader{Indicateur} &
    \tableauHeader{Cible}
}
Sauvegarder l'état civil & Fréquence & Quotidienne \\
Former les agents & Taux de formation & 100 \% en 6 mois \\
Réduire les accès partagés & Comptes nominatifs & 100 \% \\
\end{tableaucours}
\vspace{-0.8em}
\begin{center}{\small\itshape Tableau~8.2. Des objectifs de sécurité mesurables}\end{center}

\separateur

\begin{boxresume}
\textbf{Ce qu'il faut retenir de ce chapitre~:}
\begin{itemize}
    \item La \textbf{clause~4} détermine contexte, parties intéressées et
          domaine d'application.
    \item Le \textbf{domaine d'application} doit être écrit, justifié, avec ses
          exclusions.
    \item La \textbf{clause~5} exige un leadership \emph{prouvé} de la
          direction.
    \item La \textbf{clause~6} planifie les actions face aux risques et fixe
          les objectifs.
    \item Un \textbf{objectif} de sécurité est mesurable et suivi.
\end{itemize}
\end{boxresume}

\begin{boxinfo}
\textbf{Points de contrôle de l'élève.} Je sais~: déterminer le contexte et les
parties intéressées~; rédiger un domaine d'application~; citer les preuves de
leadership~; expliquer la planification~; formuler un objectif mesurable.
\end{boxinfo}

\bigskip

\begin{boxexo}[title={\ding{45}\quad Exercice~8.1. Contexte \niveaubase}]
Donnez deux enjeux externes et deux enjeux internes de la mairie d'Eséka.
\end{boxexo}

\begin{boxexo}[title={\ding{45}\quad Exercice~8.2. Parties intéressées \niveaubase}]
Citez quatre parties intéressées et une attente de sécurité pour chacune.
\end{boxexo}

\begin{boxexo}[title={\ding{45}\quad Exercice~8.3. Objectif mesurable \niveaubase}]
Transformez « mieux gérer les mots de passe » en objectif mesurable.
\end{boxexo}

\begin{boxexo}[title={\ding{45}\quad Exercice~8.4. Domaine d'application \niveaumoyen}]
Rédigez le domaine d'application du \sigle{SMSI} de la mairie, avec une
exclusion justifiée.
\end{boxexo}

\begin{boxexo}[title={\ding{45}\quad Exercice~8.5. Preuves de leadership \niveaumoyen}]
Citez quatre preuves par lesquelles le maire démontre son leadership.
\end{boxexo}

\begin{boxexo}[title={\ding{45}\quad Exercice~8.6. Périmètre flou \niveaumoyen}]
Pourquoi un domaine d'application flou rend-il l'audit impossible~?
\end{boxexo}

\begin{boxexo}[title={\ding{45}\quad Exercice~8.7. Planifier par le risque \niveaumoyen}]
Expliquez pourquoi la clause~6 renvoie à l'appréciation des risques.
\end{boxexo}

\begin{boxexo}[title={\ding{45}\quad Exercice~8.8. Trois objectifs \niveaumoyen}]
Proposez trois objectifs de sécurité mesurables pour la mairie, avec cible.
\end{boxexo}

\begin{boxexo}[title={\ding{45}\quad Exercice~8.9. Dossier de contexte \niveauexpert}]
Rédigez le dossier de contexte (clause~4) complet de la mairie d'Eséka.
\end{boxexo}

\begin{boxexo}[title={\ding{45}\quad Exercice~8.10. Plan d'objectifs \niveauexpert}]
Construisez un tableau de cinq objectifs de sécurité, avec indicateur, cible,
responsable et échéance.
\end{boxexo}

\bigskip

\begin{boxinfo}
\textbf{Pour aller plus loin.}
\begin{itemize}
    \item ISO/IEC 27001, clauses~4 à~6~\cite{iso27001}.
\end{itemize}
\end{boxinfo}

\begin{boxinfo}
\textbf{Transition.} Le Chapitre~\ref{chap:clauses710} poursuit avec le support,
le fonctionnement, l'évaluation et l'amélioration (clauses~7 à~10).
\end{boxinfo}
